% EDITED 2026-09-11 scripts/patch_koide_jordan_transcendence_2026-09-11T1930.py
% EDITED 2026-09-11 scripts/patch_koide_gap_and_parameter_count_2026-09-11T1700.py
% EDITED 2026-09-09 scripts/patch_hurwitz_citekey_2026-09-09T1640.py
%  EDITED 2026-07-12 (second pass): op:up-quark records the N -> R -> S
%  outcome -- refutation (annihilation lemma), core-coherence
%  obstruction (-3 universal), composite closure (-2 by which-path
%  recording, FANOUT), assumption ledger, and the two-item residue
%  (recording-from-FANOUT; the selection/comparability condition).
%  Classification per JS: exact algebra, circuit faces demonstrations.
%  EDITED 2026-07-12: op:up-quark brought current with the sedenion
%  results -- digit fix (+28.7% measured / +28.6% exact 9/7), notebook
%  M's route-closure recorded (bare geometry contains no 2/9; bare
%  cross-algebra vertex maximal; the dressing carries the entire
%  suppression), and the dressed-operator programme's pre-registered
%  first hypothesis stated (-2 x 1/9 = -delta_0, priced at the
%  hardware-measured leptonic ninth). Candidate status unchanged.
%  EDITED 2026-07-11: class-discreteness clarification installed after
%  the 2026-07-06 retirement of state-level GHZ/W superposition
%  (fourth_qubit doc, section 10) -- the mechanism paragraph, the
%  verified-properties paragraph (probe-family clause), and
%  op:up-quark's first clause now state explicitly that the quark
%  state remains GHZ-class and species do not superpose; the effective
%  admixture's carrier is the sedenion cross-algebra sector, as the op
%  already stated. Numbers unchanged.
%  RE-AUDITED 2026-07-06 (final pass, fourteenth stop, light leg):
%  the "previous section enumerated the spectrum" adjacency broke
%  under the reorder (confinement intervenes between the zoo and
%  here) -- retargeted to sec:generations; photon-masslessness passage
%  gains the sec:mass-gap two-way ("no debt-free direction anywhere,
%  hence a gap"); associator-higgs stitch grep EMPTY here = acquittal,
%  those seams live in interactions and boson-masses; cabibbo/sedenion
%  forward sentence confirmed present; header umlaut normalised, pure
%  ASCII. The deep audit otherwise stands.
% =====================================================================
%  Section: Masses   (Plan Stages 5-7)
%  Paper 1, "It from Bit via G\"odel". Drafted 2026-06-08.
%  Follows \S sec:generations; precedes the boson-mass / electroweak section.
%
%  \input fragment. biblatex/\autocite, amsthm `proposition'.
%  New references: koide1983new, georgijarlskog1979new
%  (delivered in masses_refs_2026-06-08.ris). J3(O) via baez2002octonions
%  and duboisviolettetodorov2019exceptional, already in the library.
%
%  Cross-reference labels (placeholders -- re-point before compiling):
%    sec:generations  -- the particle zoo / depth ordering (previous section)
%    sec:charges      -- quantum numbers; the "debt-free" photon forward-ref lands here
%    sec:register     -- the three-qubit register
%    sec:interactions -- circuits branch (forward; CKM, FCNC)
%    sec:boson-masses -- boson masses / electroweak parameters (forward)
%    sec:predictions  -- neutrino masses, glueballs (forward)
%
%  HONESTY FLAGS, now recorded as named open problems (paper-wide
%  convention; preamble in the quantum-numbers header):
%    op:associator-geodesic -- associator-debt vs geodesic proxy;
%    op:koide-delta         -- derive delta = 2/9 from Hopf holonomy
%                              (key open calculation; sin^4 holonomy route
%                              recorded as tried-and-failed);
%    op:colour-correction   -- derive |Q|^{3/2} law + phase-dilution rule;
%                              reconcile Casimir-vs-fitted Georgi-Jarlskog;
%    op:up-quark            -- the +29% discrepancy (candidate mechanism:
%                              generation-1 GHZ/W interference; enabling
%                              properties -- first-order response, phase
%                              rigidity -- verified by executed quantum
%                              computation, simulator + IBM hardware,
%                              app:qiskit-interference; r->mass map and
%                              admixture derivation still owed);
%    op:sector-scales       -- reduce the three sector scales to one;
%    op:cabibbo             -- mass phase = Cabibbo angle (two-generation
%                              form executable, passes at per-cent level:
%                              V_ud +0.11%, V_us -1.74%, simulator + IBM
%                              hardware, app:qiskit-ckm; derivation and
%                              full 3x3 CKM with CP phase still owed).
%
%  ADDED 2026-07-06: op:cabibbo gains a forward sentence naming the
%  sedenion/fourth-qubit route (development deferred; see the
%  interactions section's provenance paragraph and the session doc
%  fourth_qubit_ckm_pmns_2026-07-06.md).
%  AUDITED 2026-07-04 (joint with boson section): every claimed number
%  re-verified to the digit (script masses_boson_verification_2026-07-04.py,
%  repo scripts/); op:associator-geodesic widened to all mass per the
%  confinement unification and given its first derived constraint
%  (presentation-invariance, prop:energy-locus). No errors found.
%  REVISED 2026-06-11 after review and full numerical re-verification:
%  lepton table corrected to source values (e -0.01%, mu -0.01%; drafted
%  +0.2%/-0.03% were stale); PHASE-DILUTION RULE delta_sector = (2/9)/n
%  added to the quark subsection -- its omission made the quoted quark
%  numbers unreproducible (naive same-delta transplant fails by orders of
%  magnitude); 2/9 dimension reading corrected: dim(OP^1) = 8, not 9; the
%  correct nine-dimensional object is O+R, the ambient of the Hopf base
%  (= Bloch + colour, app:lem:base). BACKPORT: source docs carry
%  dim(OP^1) = 9 -- fix there too. One-scale claim made honest (three
%  sector scales as computed; y_t = 1 ties up-sector to v). Neutrino
%  forward-pointer: RESOLVED (2026-06-12) to ~52 meV per
%  fermion_mass_geodesic_calculation.md (dilution rule, n = 2); the
%  2.6 meV figure in quark_neutrino_masses.md is the superseded early fit
%  -- BACKPORT a supersession note there. sec:predictions (to be drafted)
%  carries: Sigma ~52 meV with the Delta-m^2_21 (x30) tension and the
%  0.6-degree theta_nu fine-tuning as a named open problem; PLUS the
%  robust qualitative set, stated as such: m_1 = 0 exactly (rank-2 seesaw
%  from two steriles), normal ordering mandatory, two sterile states
%  ~3.6 TeV, no fourth generation at any scale; falsifiers JUNO, DUNE,
%  DESI, Euclid, CMB-S4.
% =====================================================================

\section{Masses}
\label{sec:masses}

\S\ref{sec:generations} enumerated the spectrum and asserted, qualitatively, that mass grows with
the algebraic depth of the frustrated qubit. This section makes that quantitative. It begins by saying
what mass \emph{is} in this framework --- the promised definition of the ``debt'' whose absence made
the photon massless in \S\ref{sec:charges} --- and then computes two spectra: the charged leptons,
which are the clean showcase, reproduced from a single dimensional scale and no free dimensionless
parameters; and the quarks, which require two structural modifications --- one to the amplitude, one
to the phase --- and succeed for eight of the nine charged fermions, with the one failure flagged
plainly. Throughout, masses are computed on the canonical orbits of \S\ref{sec:canonical-orbits},
where the reduced spectra --- and with them the quantities below --- are sharp.

\subsection{What mass is: associator-debt}

In an associative algebra the product $(ab)c$ equals $a(bc)$ and there is nothing to record. The
octonions are not associative, and the failure is measured by the associator
\begin{equation}
  [a,b,c] = (ab)c - a(bc),
\end{equation}
which is generically non-zero. The framework identifies \emph{mass with this associator-debt}: the
information cost of specifying a bracketing that the algebra itself leaves ambiguous. Dimensionally
the associator carries an inverse length in natural units, as a mass must. A channel that carries no
associator-debt is massless --- this is exactly the ``debt-free'' direction of \S\ref{sec:charges},
and it is why the photon, the one electroweak direction with vanishing debt, has no mass. (The
colour sector's contrasting claim --- no debt-free direction anywhere, hence a gap --- is drawn at
\S\ref{sec:mass-gap}.)

It is important to be exact about what the associator does and does not fix. Computed on the basis
units, the associator has a uniform magnitude across all bracketings that are not forced to vanish;
it therefore establishes that mass exists and is non-zero, and that the three generations'
bracketings are inequivalent, but it does \emph{not} by itself set the mass hierarchy. The hierarchy
comes from how the triality structure breaks --- the angle of the next subsection --- while the
associator supplies the fact of mass. This division of labour, associator for existence and triality
angle for ratios, is the honest reading of the calculation and we keep the two separate throughout.
%% ============================================================================
%% SPLICE: the debt-to-defect bridge  (2026-09-01)
%% Target file: masses_section (current master; the 2026-06-08 lineage).
%% Suggested anchor: immediately AFTER the paragraph ending
%%   "... and we keep the two separate throughout."
%%   (the division-of-labour paragraph of \subsection{What mass is: associator-debt}).
%%   The existing "Two structural points complete the picture." paragraph then
%%   follows the splice unchanged. VERIFY the anchor whitespace-insensitively
%%   before splicing.
%% Label: eq:debt-defect -- grep the master for collisions before compiling.
%% Citation keys: harvey1982calibrated and bryant1987exceptional are NEW, carried
%%   in calibration_refs_2026-09-01.ris (web-verified 2026-09-01, DOIs included).
%%   VERIFY-KEY: replace \autocite{hurwitz1898} below with the library's existing
%%   Hurwitz key before compiling.
%% Status of claims: the displayed identity is machine-verified in convention A
%%   (sedenion_debt_koide_scan_2026-08-31.py; double-verified on both
%%   environments). The calibration identification Delta = 4(*phi), with the
%%   four-form built independently via the Hodge star, is verified in
%%   defect_calibration_scan_2026-09-01.py (single environment; echo owed).
%%   Orientation sign sigma = +1 in convention A.
%% House CAUTION honoured: "bump in the carpet" below refers to the algebraic
%%   instance only (the computability split), per the instance-discipline
%%   comment in the computability-split introduction.
%% ============================================================================
%% BEGIN SPLICE bridge-debt-defect
A word on the provenance of the term, because its standing has changed.
\emph{Associator-debt} began as a coinage of this programme: we conceived of it
as a logical debt --- the information owed by a bracketing that the algebra
never discharges, the bump in the carpet in its algebraic instance --- and we
kept it because it priced the mass sector correctly, not because the object had
an address in the wider literature. It acquired one when we stopped treating
the sedenion zero divisors as a bare obstruction and asked what structure they
carry. At the sedenion step, exactly where Hurwitz's theorem ends the division
tower \autocite{hurwitz1898composition}, the norm of a product first fails to factorise,
and the failure \emph{is} the debt:
\begin{equation}
  \bigl\lvert (a,b)(c,d) \bigr\rvert^{2}
  - \bigl\lvert (a,b) \bigr\rvert^{2} \bigl\lvert (c,d) \bigr\rvert^{2}
  \;=\; 2 \,\bigl\langle a,\, [c,b,d] \bigr\rangle ,
  \label{eq:debt-defect}
\end{equation}
and the right-hand side is, up to normalisation, the coassociative calibration
four-form of the $G_2$-structure evaluated on the four octonionic components
\autocite{harvey1982calibrated,bryant1987exceptional}, with the zero divisors
sitting exactly at the calibration bound. \emph{Norm defect} is therefore
another way of saying \emph{associator-debt}: the associator prices the debt at
issuance, where the octonion norm still balances the books; the sedenion norm
defect is where it settles; and the settlement geometry is standard calibrated
geometry rather than a private idiom. The identity fixes the object, not the
spectrum --- the division of labour above stands, with the generation splitting
still supplied by the frame and the $\mathbb{Z}_3$ structure of the tag.
%% END SPLICE bridge-debt-defect

Two structural points complete the picture. First, because the non-associativity forbids a metric on
the octonionic fibre $S^7$, mass cannot literally be a geodesic length on that fibre; it is the
associator-debt. The operational calculations below, which are sometimes phrased as ``geodesic''
quantities, are a computational proxy.

\begin{openproblem}[Associator-debt and the geodesic proxy]\label{op:associator-geodesic}
Show that the geodesic-length proxy used in the operational mass calculations computes the same
quantity as the associator-debt that defines mass, closing the gap between definition and
computation. The reconciliation must hold for the whole mass sector: the fermion spectra of this
section and, by the unification of \S\ref{sec:confinement}, the glueball and binding networks as
well. And one constraint on any solution is already a theorem: on closed networks the map must be
presentation-invariant, attaching to the network as such and to no particular cut
(Proposition~\ref{prop:energy-locus}).
\end{openproblem}

Second, the dimensional-parameter count. Structurally the framework requires exactly one free
dimensional parameter --- an overall scale --- and this is a necessity rather than a gap: the
observer cannot fix the absolute scale of their own computational network from inside it
(\S\ref{sec:register}). The Higgs vacuum value sets that scale, and the doubling norm $\sqrt2$ that
appears below is the algebraic signature of the vacuum value ``doubling'' the structure. As the
calculation currently stands, however, each charge sector carries its own fitted scale --- three in
all --- of which the up-quark scale is tied to the electroweak value by $y_t = 1$
(\S\ref{sec:boson-masses}), while the ratio of the down-quark to the lepton scale sits suggestively
near the doubling norm squared, $M^2_{\mathrm{down}}/M^2_{\mathrm{lepton}} \approx 2.07$. Reducing
the three to the structurally required one is open.

\begin{openproblem}[The sector scales]\label{op:sector-scales}
Derive the lepton and down-quark sector mass scales from the single electroweak scale, as
$y_t = 1$ already does for the up-quark sector --- or show why the residual ratios
($\approx 2$ for down-to-lepton) are fixed by the algebra.
\end{openproblem}

\subsection{Charged leptons: the Koide relation}

The charged-lepton masses satisfy the empirical relation of Koide \autocite{koide1983new},
\begin{equation}
  \frac{m_e + m_\mu + m_\tau}{\bigl(\sqrt{m_e}+\sqrt{m_\mu}+\sqrt{m_\tau}\bigr)^2} = \frac{2}{3},
  \label{eq:koide}
\end{equation}
which holds to about one part in $10^5$ and has resisted explanation for forty years. The framework
produces it as follows. The three generations are the three off-diagonal octonionic entries of the
exceptional Jordan algebra $J_3(\mathbb{O})$ \autocite{baez2002octonions,duboisviolettetodorov2019exceptional},
cyclically permuted by a $\mathbb{Z}_3$ triality, so the square-root masses are the three triality
images of one quantity,
\begin{equation}
  \sqrt{m_k} = M\bigl(1 + \alpha\cos(\delta + \tfrac{2\pi k}{3})\bigr), \qquad k = 0,1,2,
  \label{eq:z3}
\end{equation}
with $M$ a scale, $\alpha$ the amplitude, and $\delta$ a phase. The amplitude is the Cayley--Dickson
doubling norm $\alpha = \sqrt2$ ($|1+i| = \sqrt2$), and this single value forces the Koide ratio.

\begin{proposition}[The Koide ratio from the doubling norm]\label{prop:koide}
For square-root masses of the form~\eqref{eq:z3}, the ratio~\eqref{eq:koide} equals
$(1 + \tfrac12\alpha^2)/3$, independently of $M$ and $\delta$. With $\alpha^2 = 2$ it equals $2/3$.
\end{proposition}

\begin{proof}
The three phases are spaced by $2\pi/3$, so $\sum_k \cos(\delta + 2\pi k/3) = 0$ and
$\sum_k \cos^2(\delta + 2\pi k/3) = \tfrac32$. Then $\sum_k\sqrt{m_k} = 3M$, so the denominator is
$9M^2$, while $\sum_k m_k = M^2\sum_k(1+\alpha\cos)^2 = M^2(3 + \tfrac32\alpha^2) = 3M^2(1+\tfrac12\alpha^2)$.
The ratio is $(1+\tfrac12\alpha^2)/3$, and $\alpha^2=2$ gives $2/3$.
\end{proof}

The Koide relation is thus derived, not fitted: it is the statement that the triality amplitude is
the doubling norm. The phase is $\delta = 2/9$ in radians, the value that the measured masses
require, matched to five significant figures and carrying a natural geometric reading as
$2/9 = \dim_{\mathbb{R}}\mathbb{C} \,/\, \dim_{\mathbb{R}}(\mathbb{O}\oplus\mathbb{R})$: the
preferred complex plane against the nine-dimensional ambient space of the Hopf base $S^8$ ---
exactly the (Bloch)$\,\oplus\,$(colour) space of Appendix~\ref{app:setup}. The reading is
suggestive, matched rather than derived. With $\alpha$ fixed by the theory, $\delta$ matched once,
and the single scale $M$ set by the heaviest mass, the other two follow:

\begin{center}
\begin{tabular}{lccc}
\hline
lepton & predicted (MeV) & observed (MeV) & error \\
\hline
$e$    & $0.5110$  & $0.5110$  & $-0.01\%$ \\
$\mu$  & $105.65$  & $105.66$  & $-0.01\%$ \\
$\tau$ & $1776.9$  & $1776.9$  & (scale) \\
\hline
\end{tabular}
\end{center}

This is the showcase: three masses from one dimensional input and no fitted dimensionless
parameters, the dimensionless content being the two numbers $\alpha=\sqrt2$, which the theory
fixes, and $\delta=2/9$, which it matches.

We flag the one genuine gap squarely. Proposition~\ref{prop:koide} derives the $2/3$ relation from
$\alpha=\sqrt2$ and the $\mathbb{Z}_3$ ansatz, both of which the framework supplies; but the specific
phase $\delta = 2/9$ is at present a \emph{matched} value with a suggestive geometric interpretation,
not a derived one.

\begin{openproblem}[The Koide phase]\label{op:koide-delta}
Derive $\delta = 2/9$ from the holonomy of the octonionic Hopf connection. This is the single
most important open calculation in the mass sector. A derivation passes if it lands inside the
measured band, $\delta = 0.222225 \pm 0.000006$ radians, with no fitted input; the invariant
target is $3\delta = 2/3$, the phase of the cyclic product of the three off-diagonal units
(Proposition~\ref{prop:koide-jordan}). It fails before any number is checked if it ends in an
algebraic value of $e^{i\delta}$ --- a root of unity, the argument of a configuration of vectors
with algebraic coordinates, the phase of a cyclic product of algebraic octonions --- since
$e^{2i/9}$ is transcendental (\S\ref{sec:koide-gap}). Four routes are closed: the simplest
holonomy-squared model, $m_k \propto \sin^4(\theta_0 + 2\pi k/3)$, does not reproduce the
spectrum; the structural angles of the bare sedenion triple contain no $2/9$
(Appendix~\ref{app:nb-M}); the norm-defect identity~\eqref{eq:debt-defect}, being a magnitude,
cannot fix a phase; and a coherent-state realisation of the triality frame, with the phase as
the enclosed phase-space area, cannot meet the amplitude and the phase at once.
\end{openproblem}

\subsection{The shape of the gap}
\label{sec:koide-gap}

Since the phase is the one number in the charged-lepton sector that the framework matches rather
than derives, it is worth being exact about what kind of gap it is, what has been tried, and
what a derivation would have to look like. A well-posed gap is an invitation; a vague one is a
hole.

What the data fix. Inverting the triality form~\eqref{eq:z3} on the three measured masses is
exact --- three masses, three parameters --- and with the PDG uncertainties propagated it
returns $\alpha = 1.41421 \pm 0.00001$, within half a standard deviation of $\sqrt2$, and
$\delta = 0.222225 \pm 0.000006$ radians, within half a standard deviation of $2/9 = 0.222222$.
The band is set almost entirely by the tau mass and fixes between four and five significant
figures of $\delta$; any candidate derivation is tested against that band and nothing looser.

Why the norm defect does not close it. The identity~\eqref{eq:debt-defect} is quartic in the
magnitudes: it says how badly the norm of a product fails to factorise and bounds the failure by
the coassociative calibration. Everything it can fix is a size --- the existence of mass, the
scale of the debt, the amplitude $\alpha$, which is itself a norm. The phase $\delta$ is not a
size but an orientation, the angle of the $\mathbb{Z}_3$ triality frame against the preferred
complex structure, and a magnitude identity has no handle on an orientation. That is the content
of the sentence that closes the bridge above: the identity fixes the object, not the spectrum.
The division of labour --- existence from the associator, ratios from the frame --- holds at the
sedenion rung exactly as at the octonion rung, and the elegance of the norm defect as an account
of mass does not carry over to the phase.

What is already excluded. Three routes are closed. The simplest holonomy-squared model fails
(Open Problem~\ref{op:koide-delta}). The bare sedenion geometry was inventoried in the up-quark
investigation (Appendix~\ref{app:nb-M}): an eighteen-comparison scan of the structural angles of
the sedenion triple --- twelve identically zero, four at exactly $\pi/4$, two ratios at exactly
$\tfrac14$ --- contains no occurrence of $2/9$. That scan was run against the up-quark leakage
angle, but it is the same number and the same geometry, and it closes the frame-orientation
route for the Koide phase as well. And the associator has a uniform magnitude on every
bracketing not forced to vanish, by $G_2$-transitivity, so no choice of octonionic direction can
supply the number either.

The shape any derivation must take. The number $2/9$ is a rational number of \emph{radians}, and
that is an unusual thing for a geometric angle to be. An angle fixed by a discrete symmetry is a
rational multiple of $\pi$; the inventory's $\pi/4$ is what such angles look like. An angle
between two vectors is an inverse trigonometric function of an algebraic number. A holonomy on a
round sphere around a loop singled out by symmetry is a rational fraction of $4\pi$, and a loop
not singled out by symmetry is a fitted input by another name. A pre-registered scan of these
families against the measured band finds no member of any of them at a simple denominator: the
nearest rational multiple of $\pi$ with denominator up to sixty lies nearly three hundred
standard deviations from the band, and the nearest inverse trigonometric value of a rational or
a surd with denominator up to a hundred lies ten away. (Pushing the denominators to four hundred
eventually produces a chance coincidence, $22\pi/311$, at exactly the rate a random search
would; it carries no information.) The dimension-ratio reading $2/9 =
\dim_{\mathbb{R}}\mathbb{C} \,/\, \dim_{\mathbb{R}}(\mathbb{O}\oplus\mathbb{R})$ meets the same
difficulty in another form. A ratio of dimensions is a pure number, and turning it into an angle
requires a unit; the natural unit of angle in a symmetric geometry is the turn, and two-ninths
of a turn is $1.396$ radians, nowhere near. The reading needs a mechanism in which one radian is
the natural unit --- in which the angle is an arc measured against its radius, two units of arc
on a radius of nine --- and until that mechanism is exhibited the reading remains what we called
it: suggestive, matched rather than derived.

What the phase is. The triality form~\eqref{eq:z3} is not an ansatz laid over the Jordan
algebra; it is the spectrum of the simplest element of it.

\begin{proposition}[The triality form as the spectrum of a democratic Jordan element]\label{prop:koide-jordan}
Let $X \in J_3(\mathbb{O})$ have unit diagonal and off-diagonal entries $x_1, x_2, x_3$ of a
common norm $\rho$, and write $x_i = \rho u_i$ with $u_i$ unit octonions and
$\operatorname{Re}(u_1u_2u_3) = \cos\Phi$. Then the eigenvalues of $X$ are $1 +
2\rho\cos\bigl((\Phi + 2\pi k)/3\bigr)$, $k = 0, 1, 2$: the form~\eqref{eq:z3} with $\alpha =
2\rho$ and $3\delta = \Phi$. The phase is well defined although the octonions are not
associative, and for $\rho = 1/\sqrt2$ it satisfies $|\Phi| \le \pi/4$, with equality exactly
when one eigenvalue vanishes.
\end{proposition}

\begin{proof}
The characteristic cubic of $X$ is $\lambda^3 - (\operatorname{tr}X)\lambda^2 + S(X)\lambda -
N(X)$ with $\operatorname{tr}X = 3$, $S(X) = 3 - 3\rho^2$ and cubic norm $N(X) = 1 - 3\rho^2 +
2\rho^3\cos\Phi$ \autocite{baez2002octonions}. Substituting $\lambda = 1 + 2\rho c$ reduces it
to $4c^3 - 3c = \cos\Phi$, whose roots are $c = \cos\bigl((\Phi + 2\pi k)/3\bigr)$. The real
part of a triple product is the same for either bracketing in any composition algebra, which is
what makes $\Phi$ well defined. For $\rho = 1/\sqrt2$ the smallest eigenvalue is $1 +
\sqrt2\cos\bigl((|\Phi| + 2\pi)/3\bigr)$, which is non-negative exactly when $|\Phi| \le \pi/4$.
\end{proof}

The proposition says what the two numbers are inside the framework's own construction. The
amplitude $\alpha = \sqrt2$ is the statement that the off-diagonal entries have squared norm one
half; read through the Gram matrix of three generation states, whose entries are the overlaps
$\langle\psi_j|\psi_k\rangle$, it says that each pair is mutually unbiased,
$|\langle\psi_j|\psi_k\rangle|^2 = \tfrac12$, and the Koide ratio then follows exactly as in
Proposition~\ref{prop:koide}. The phase is the Bargmann invariant of the triple, the argument of
the cyclic product of the three overlaps, which is the one continuous invariant three states
possess beyond their overlaps; its octonionic form is the real part of the cyclic product of the
three off-diagonal units. The measured $3\delta = 0.6667$ sits at $0.85$ of the positivity bound
$\pi/4$, and to first order the electron's square-root mass is one third of the shortfall,
$\sqrt{m_e}/M \simeq (\pi/4 - 3\delta)/3$, accurate to two per cent. So the question ``why
$2/9$'' is, precisely, ``why does the cyclic product of the three off-diagonal units have real
part $\cos(2/3)$''. The number to derive is the $\mathbb{Z}_3$-invariant $3\delta = 2/3$, and it
is worth recording that this is also the value of the Koide ratio: the invariant phase, in
radians, equals $\operatorname{tr}X^2/(\operatorname{tr}X)^2$. We have no mechanism that makes a
phase equal a trace ratio and record the equality as a rhyme.

What kind of number it must be. If $\delta = 2/9$ exactly then $e^{i\delta}$ is transcendental,
by the Lindemann--Weierstrass theorem: $e^{a}$ is transcendental for every non-zero algebraic
$a$. The consequences are sharp. The cyclic product $u_1u_2u_3$ cannot be an algebraic octonion,
so no construction of the three entries from the structure constants, from roots of unity, from
a finite group, or from any configuration of vectors with algebraic coordinates can be exact;
the holonomy of a round connection around a geodesic polygon with algebraic vertices is the
argument of exactly such a product and is excluded with them. Every charged-lepton mass ratio is
then a transcendental number, since each is a non-constant rational function of $e^{i\delta}$
with algebraic coefficients. Conversely, any mechanism that returns an algebraic value for a
single mass ratio returns a $\delta$ that is zero or transcendental, and must then explain the
agreement with $2/9$ as an accident: a fraction with denominator at most nine lands in the
six-sigma band by chance about once in a thousand trials. What survives is the class in which
the phase is the exponential of a rational number --- an action or a symplectic area in units of
$\hbar$, a length along a flat direction in units of its radius, a Lie-algebra generator with a
rational coefficient. That is the precise content of ``flat rather than round''.

Flat, but not a torus. On a compact flat torus a quantised flux makes every enclosed phase a
rational multiple of $2\pi$, which is the family the scan excludes; and a flat connection whose
holonomy is not quantised is a modulus, a free parameter that the geometry leaves open. The flat
structure carrying the phase must be non-compact in the direction that matters, and the rational
number is then most naturally an area, the product of two rational lengths in units of $\hbar$.
The one such reading we have found is a rhyme: $3\delta = 2/3 = 2\cdot(1/\sqrt3)^2$ is the
commutator phase of two orthogonal phase-space displacements of squared amplitude one third, the
equal-weight number. The obvious model built on it --- three coherent states at the vertices of
an equilateral triangle, mass roots the eigenvalues of their Gram matrix, phase twice the
enclosed area --- has one parameter and two targets and fails both ways: fixing the overlap at
$1/\sqrt2$ gives a phase of $0.600$ rather than $0.667$, and fixing the phase gives an amplitude
of $1.361$ rather than $\sqrt2$. We record the shape of the derivation as a constraint, the
model as a closed route, and the earlier suggestion of a torus as superseded.

\subsection{Quarks: the colour correction}

Quarks do not satisfy the Koide relation --- the ratio~\eqref{eq:koide} is about $0.73$ for the
down-type and $0.85$ for the up-type quarks rather than $2/3$, a deviation that is invariant under
perturbative QCD running and therefore a bare structural fact, not a scale artefact. The framework
reads it as the effect of the colour debt that leptons, being colour singlets, do not carry. Two
modifications enter, one to the amplitude and one to the phase. Where the colourless leptons sit at
the bare doubling norm $\alpha^2 = 2$, the coloured quarks acquire an additional contribution that
scales with electric charge,
\begin{equation}
  \alpha^2 = 2 + 2\,|Q|^{3/2},
  \label{eq:quarkalpha}
\end{equation}
the coefficient being the doubling norm squared and the power $3/2$ fixed empirically to within one
per cent; the framework reads $|Q|^{3/2} = |Q|\cdot\sqrt{|Q|}$ as winding number times its
amplitude --- the Born-rule pairing appearing inside the fibre metric --- but this reading is a
gloss, not a derivation. The formula gives $\alpha^2 = 2.385$ for the down-type quarks
($|Q|=\tfrac13$) and $3.089$ for the up-type ($|Q|=\tfrac23$), against $2.389$ and $3.094$ extracted
from the measured sector ratios. The phase is \emph{diluted}: each sector's triality phase is
\begin{equation}
  \delta_{\mathrm{sector}} = \frac{2/9}{n},
  \qquad
  n \;=\; 1 + \bigl(T_3 + \tfrac12\bigr) +
  \begin{cases}1 & \text{colour triplet}\\[2pt] 0 & \text{colour singlet,}\end{cases}
  \label{eq:dilution}
\end{equation}
where $n$ counts the active fibre directions the channel engages --- the $S^1$ always, the $S^3$
for upper isospin, the $S^7$ for colour --- among which the misalignment $2/9$ is shared. Thus
$n = 1$ for the charged leptons ($\delta = 2/9$, as above), $n = 2$ for the down-type quarks
($\delta = \tfrac19$), and $n = 3$ for the up-type ($\delta = \tfrac{2}{27}$); triality shifts of
$2\pi/3$ relabel the generations without changing the spectrum.
Feeding~\eqref{eq:quarkalpha} and~\eqref{eq:dilution} into the triality form~\eqref{eq:z3}, with the
scale set by the heaviest quark in each sector, the lighter masses follow: $d$ at $4.62$~MeV
($-1.0\%$) and $s$ at $94.4$~MeV ($+1.1\%$) in the down sector; $c$ at $1271$~MeV ($+0.1\%$) in the
up sector. Eight of the nine charged fermions are reproduced to about one per cent.

The ninth is not. The up quark comes out at $2.78$~MeV against an observed $2.16 \pm 0.07$~MeV, an
error of some twenty-nine per cent and far outside the quoted uncertainty. This is a genuine failure
of the formula as it stands, not a rounding issue. Two mitigating
observations follow, neither a resolution. The up-quark mass is never measured directly --- it is
extracted from lattice QCD with chiral perturbation theory, and determinations range over roughly
$1.9$--$2.5$~MeV, at the top of which range the discrepancy shrinks to about twelve per cent. And a
candidate mechanism is identified: at the first generation, and only there, the tripartite (GHZ) and
pairwise (W) contributions to the mass are of comparable size, and their interference shortens the
effective path. The clause that keeps this consistent with the framework's species assignment: the
quark state remains GHZ-class throughout --- the pairwise structure is forced \emph{within} that
class by the Cayley--Dickson nesting $\mathbb{C} \subset \mathbb{H} \subset \mathbb{O}$, and
the interference is between contributions to the mass functional, never a superposition of
species, which are class-discrete. The mechanism is consistent with the up quark also being the
one fermion whose $T_3 = +\tfrac12$ member is \emph{lighter} than its doublet partner, the
competition between the amplitude effect (larger $\alpha$ for up-type) and the dilution effect
(smaller $\delta$) tipping the other way at the lightest generation.

The candidate mechanism's enabling properties have since been verified by an executed quantum
computation (Appendix~\ref{app:qiskit-interference}) on a one-parameter probe family that
traverses the GHZ class toward the W boundary --- the class-discreteness clarification is stated
there --- run on a statevector simulator and subsequently on IBM Quantum hardware. Three facts are established. The interference exists in
exactly the invariant the geodesic construction is built from: the Bloch radius of the cut responds
to a GHZ--W admixture. The response is \emph{first order} in the admixture amplitude --- a
two-per-cent W component moves the invariant seventeen times more than incoherent mixing would ---
so a modest admixture is large enough to produce an effect of the required size. And the response
is rigid in the relative phase, so the mechanism, if it is the right one, must live in the
admixture magnitude. What the computation does not yet supply is the map from the invariant to the
mass --- the reconciliation of Open Problem~\ref{op:associator-geodesic} --- nor the
framework-derived admixture for the first generation, nor, therefore, the sign and size of the
resulting shift, which must come out as a reduction of the predicted value by some twenty per
cent. The hypothesis is thus verified as viable and sharply constrained, not as confirmed; the
distinction is exactly the one this paper's conventions exist to keep. The required reduction has
since acquired a candidate identity: it is $\delta_0 = \tfrac29$ itself ---
$m_u = (1 - \tfrac29) \times 2.78 = 2.162$~MeV against the observed $2.16$ --- the Koide angle in
its fourth role, linear in the interference parameter exactly as the verified first-order response
requires. The identity was found by a disclosed candidate scan and is recorded, not derived; the
experimental band on $m_u$ is far wider than the central-value agreement suggests, and the
open problem below states what deriving it would take.

\begin{openproblem}[The up quark]\label{op:up-quark}
Give the generation-one interference mechanism a formal treatment: derive the effective
admixture magnitude from the framework (the quark state remaining GHZ-class throughout; species
do not superpose), supply the map from the cut's Bloch radius to the mass
(Open Problem~\ref{op:associator-geodesic}), and determine whether the resulting shift accounts,
in sign and size, for the up-quark discrepancy ($2.78/2.16 = +28.7\%$ as measured; exactly
$\tfrac{9}{7} - 1 = +28.6\%$ if the candidate identity below holds) --- or find what does. The mechanism's
enabling properties, the first-order response and the phase rigidity verified in
Appendix~\ref{app:qiskit-interference}, are constraints any formal treatment must respect. The
problem now has a sharp form: a recorded candidate takes the admixture magnitude to be the Koide
angle itself, $m_u = (1 - \delta_0) \times$ the formula value, landing on the observed mass at
central-value precision with no new constant --- derive why the generation-one leakage angle
equals $\delta_0$, with the sedenion cross-algebra sector (\S\ref{sec:sedenion-announcement})
the candidate home of the computation, or refute the identity. The bare geometry has since been
scanned and does not supply the angle: an eighteen-comparison inventory of the sedenion triple's
structural angles contains no occurrence of $\tfrac29$, and the bare cross-algebra vertex mixes
\emph{maximally} --- so the identity, if it holds, is a property of the \emph{dressed} leakage
operator, and the suppression from maximal to $\delta_0$ is the object the derivation must
produce, not a constant it may look up. The programme's pre-registered first hypothesis, its pass
criteria fixed before any computation: the first-order dressed shift is
$-2 \times \tfrac19 = -\delta_0$ --- two elementary cross-algebra moves, each priced at the
leptonic ninth of Appendix~\ref{app:hw-ninth} --- the two-ninths identity arriving at the mass
spectrum.
% CUTOVER: add \ref{app:nb-M} for the eighteen-comparison scan when the notebook record is input.
The programme has since run, in three registered stages. The first family --- the dressing as
the raw composition-failure operator --- was refuted by a two-line lemma: the defect operator
annihilates the very generation line whose failure it encodes. The second established the
core-coherence obstruction: every zero-divisor excursion decomposes half into the shared
quaternionic core and half into a single third-sector unit, the core halves add coherently, the
coefficient is $-3$ for every diagonal coupling, and the required $-2$ is unreachable on the
algebra alone. The third closed the pricing on the register--algebra composite: each excursion
carries its generation tag, the which-path record decoheres the core coherence ---
\textsc{fanout} doing its defining job --- and the coefficient is exactly $-2$, so
$\delta m / m = -2 \times \tfrac19 = -\delta_0$ with the assumption ledger reduced to named
items: channel count two (derived, the Klein grading); the $S_3$ twist (derived); the price
(the measured ninth); excursions register-recorded (the composite's defining property). Two
derivation targets remain and are this problem's residue: derive the recording clause from the
\textsc{fanout} boundary rather than assume it; and derive the \emph{selection} --- the
dressed shift is generation-universal and the induced mixing vanishes identically (an exact
cancellation arranged by the twist), so the identity's visibility at the first generation alone
is carried by the comparability condition above, now cleanly separated from the pricing. The
computations behind these statements are exact sedenion linear algebra, reproducible without
quantum software; their circuit faces are demonstrations, of which the which-path pair --- the
obstruction's $3$ against the record's $2$ --- is the one a physical device adds meaning to.
% CUTOVER: add \ref{app:nb-N}, \ref{app:nb-R}, \ref{app:nb-S} when the notebook record is input.
\end{openproblem}

There is a further structural thread that we flag rather than resolve. The quark--lepton mass
relations carry a factor of three of the Georgi--Jarlskog type \autocite{georgijarlskog1979new},
which the framework can read either group-theoretically, as a Casimir ratio of the colour structure,
or through the fitted charge dependence of~\eqref{eq:quarkalpha}. What the framework does \emph{not}
need is the $SU(5)$ scaffolding through which the Georgi--Jarlskog factor is conventionally obtained
--- consistent with the exclusion of $SU(5)$ and the absence of proton decay noted in
\S\ref{sec:charges}.

\begin{openproblem}[The colour correction]\label{op:colour-correction}
Derive the amplitude law~\eqref{eq:quarkalpha} --- in particular the power $3/2$ --- and the
phase-dilution rule~\eqref{eq:dilution} from the octonionic structure of the colour-triplet Higgs
coupling, and reconcile the Casimir and fitted accounts of the Georgi--Jarlskog factor.
\end{openproblem}

\subsection{Mass eigenstates versus weak eigenstates}

The basis that diagonalises the associator-debt --- the mass basis --- is not the basis in which the
weak interaction acts, which is the doublet basis of the $B$-qubit (\S\ref{sec:charges}). Their
mismatch is the Cabibbo--Kobayashi--Maskawa matrix. The framework ties the rotation to the same
octonionic geometry that fixes the masses: the lepton Koide phase $\delta = 2/9$ radians is about
$12.7^\circ$, strikingly close to the Cabibbo angle, which suggests that one geometry sets both the
mass hierarchy and the generation mixing. We mark that identification as suggestive rather than
derived --- but it is no longer merely an observation about two numbers. The identification has
been made executable (Appendix~\ref{app:qiskit-ckm}): a charged-current circuit whose $W$ gate
carries the rotation $R_y(2\delta)$ returns $\lvert V_{ud}\rvert = \cos(2/9) = 0.97541$ against
the measured $0.97435$ ($+0.11\%$) and $\lvert V_{us}\rvert = \sin(2/9) = 0.22040$ against
$0.22430$ ($-1.74\%$), run on a statevector simulator and reproduced on IBM Quantum hardware to
within shot noise. The two-generation block thus passes its first quantitative test at the
per-cent level; what is suggestive is now also \emph{checkable}, by anyone, in one rotation angle.

\begin{openproblem}[The Cabibbo angle]\label{op:cabibbo}
Derive, rather than observe, the relation between the mass phase and the generation mixing: whether
the Cabibbo angle equals the undiluted phase $\delta = 2/9$, and what fixes the remaining CKM
structure. The two-generation form is executable and passes at the per-cent level
(Appendix~\ref{app:qiskit-ckm}); the derivation, the residual $-1.74\%$ on
$\lvert V_{us}\rvert$, and the extension to the full three-generation matrix with its
CP-violating phase are what remain. A candidate route to the full matrix is recorded in
\S\ref{sec:interactions}: the generation wire read as the sedenion doubling, with mixing forced
by the four-qubit entanglement obstruction; its development is deferred to a dedicated paper.
\end{openproblem}

Its structural consequence, however, is exactly what the interaction picture needs: the mass
eigenstates are distinct from the weak eigenstates, so a charged-current vertex changes flavour while
the neutral currents, acting on the generation-blind structure, do not. That is the seed of the
no-tree-level flavour-changing-neutral-current result, which is delivered with the circuit picture of
interactions in \S\ref{sec:interactions}.

\subsection{What is settled, and what is owed}

Settled: mass is associator-debt, with the photon massless because it is debt-free; the Koide $2/3$
relation follows from the doubling norm $\alpha=\sqrt2$ and the $\mathbb{Z}_3$ triality of $J_3(\mathbb{O})$
(Proposition~\ref{prop:koide}); the charged-lepton spectrum is reproduced from one dimensional scale
and the two dimensionless numbers $\alpha=\sqrt2$, $\delta=2/9$, the electron and the muon each to
$0.01\%$; and the quark colour correction~\eqref{eq:quarkalpha}, with the phase-dilution
rule~\eqref{eq:dilution}, reproduces eight of nine charged fermion masses to about a per cent.

Owed, and recorded as the named problems above: the derivation of $\delta = 2/9$ from the Hopf
holonomy (Open Problem~\ref{op:koide-delta}, the key open calculation); the associator--geodesic
reconciliation (Open Problem~\ref{op:associator-geodesic}); the up-quark discrepancy
(Open Problem~\ref{op:up-quark}); the derivation of the colour correction and the dilution rule
(Open Problem~\ref{op:colour-correction}); the reduction of the sector scales
(Open Problem~\ref{op:sector-scales}); and the Cabibbo identification
(Open Problem~\ref{op:cabibbo}).

With the fermion masses in hand, two threads remain. The electroweak gauge bosons, whose masses use
the weak mixing angle $\sin^2\theta_W = \tfrac14$ already derived in \S\ref{sec:charges}, are taken
up in \S\ref{sec:boson-masses}; and the neutrino masses, which follow from the sterile seesaw of
\S\ref{sec:generations} and the dilution rule~\eqref{eq:dilution} with $n = 2$ --- a massless
lightest state, normal ordering, and a sum near $52$~meV --- are collected with the other
falsifiable predictions in \S\ref{sec:predictions}, where the one tension in the neutrino fit (the
first-to-second-generation splitting) is recorded as open.
